\documentclass[11pt]{article}
\usepackage[margin=1in]{geometry}
\usepackage{amsmath,amssymb,amsthm,booktabs}
\usepackage[hidelinks]{hyperref}
\setlength{\emergencystretch}{2em}
\newtheorem{theorem}{Theorem}
\newtheorem{proposition}[theorem]{Proposition}
\newtheorem{corollary}[theorem]{Corollary}
\newcommand{\mr}{\operatorname{mr}}
\newcommand{\rk}{\operatorname{rank}}
\newcommand{\Epi}{\operatorname{Epi}}
\newcommand{\Hom}{\operatorname{Hom}}
\newcommand{\Aut}{\operatorname{Aut}}
\title{Joint simple quotients of branched-cover groups}
\author{Internal research working proof}
\date{16 November 2026}
\begin{document}
\maketitle

The user's proposed test is correct: twenty distinct $A_5$ quotient
kernels of the \emph{same} double-branched-cover group imply
$\mr(K)\ge4$. This combines the retained cover inequality with a
classical finite-group threshold. Burness--Crestani explicitly give
$\delta(A_5)=19$ and credit the direct-power formula to Philip Hall
\cite[p.~2]{bc}. The application was proposed by the user; no priority
claim is made for the classical group theory or the resulting knot bound.

An exact search on all 293 names in the pinned uncertainty panel gives
at most five kernels on any one knot. A known positive control,
$P(2,3,3,5)$, has 23 kernels: the test proves $\mr=b=4$ there, whereas
the raw $A_5$ homomorphism count gives only $\mr\ge3$. The control
belongs to the already-settled Montesinos class. No previously
unresolved knot is settled, and full MRC remains open.

\section{The same-group kernel bound}

For a finite group $S$ and finitely generated group $H$, put
\[
 \kappa_S(H)=\#\{N\triangleleft H:H/N\cong S\},\qquad
 h_d(S)=\frac{|\Epi(F_d,S)|}{|\Aut(S)|}.
\]
An epimorphism is counted once per kernel, rather than once per
labelling of its target.

\begin{theorem}\label{thm:capacity}
If $\rk H\le d$, then $\kappa_S(H)\le h_d(S)$. Consequently, for
the simple degree-$r$ branched cover $Y_\phi$ associated to a
transposition quotient of a knot group,
\[
 \kappa_S(\pi_1Y_\phi)>h_d(S)\quad\Longrightarrow\quad
 \mr(K)\ge r+d.
\]
\end{theorem}
\begin{proof}
Two epimorphisms $H\to S$ have the same kernel exactly when they differ
by postcomposition with an automorphism of $S$. The automorphism action
on epimorphisms is free, so $\kappa_S(H)=|\Epi(H,S)|/|\Aut(S)|$.
If $F_d\twoheadrightarrow H$, precomposition injects the automorphism
orbits for $H$ into those for $F_d$. This proves the capacity bound.
Exceeding it forces $\rk H\ge d+1$. Apply the retained uniform
inequality $\mr(K)\ge r-1+\rk\pi_1Y_\phi$ from the companion
branched-cover note. No rank--Heegaard-genus equality is assumed.
\end{proof}

For $A_5$, there are 2280 ordered generating pairs and 120
automorphisms. Thus $h_2(A_5)=19$.
\begin{corollary}\label{cor:twenty}
If $H=\pi_1(\Sigma_2K)$ has at least twenty distinct $A_5$ quotient
kernels, then $\rk H\ge3$ and $\mr(K)\ge4$. A certified four-bridge
presentation then proves $\mr(K)=b(K)=4$.
\end{corollary}

\section{Why the small quotients act jointly}

For nonabelian simple $S$, distinct kernels have an equivalent
direct-product interpretation.
\begin{proposition}
If $f_1,\ldots,f_t:H\twoheadrightarrow S$ have pairwise distinct
kernels, their joint map $H\to S^t$ is surjective.
\end{proposition}
\begin{proof}
Induct on $t$. The image is a subdirect product of $S^{t-1}$ and $S$.
By Goursat's lemma, a proper such subgroup identifies the last factor
with a simple quotient of $S^{t-1}$. Every epimorphism
$S^{t-1}\to S$ depends on exactly one coordinate: the coordinate images
are commuting normal subgroups of the simple target, so one is all of
$S$ and all others are trivial. The identification would therefore
make $f_t$ an automorphism followed by one earlier coordinate map,
contradicting distinct kernels. Hence the image is all of $S^t$.
\end{proof}

In particular, twenty kernels give a quotient $A_5^{20}$, which cannot
be two-generated. The elementary capacity proof above already suffices
for the rank bound. Maps from twelve different knot covers cannot
define this joint map: they have different domains. Maps obtained
from several presentations require explicit transport to one fixed
group before their kernels can be compared.

\paragraph{Exact finite capacities.}
The retained permutation model lists the sixty even permutations of
five letters. Exhaustive closure gives its 59 subgroups. It has exactly
five subgroups of order twelve, and conjugation on these is faithful.
Every automorphism permutes those five subgroups. If it fixes all five,
faithfulness of the conjugation action forces it to fix every element
of $A_5$. Thus $\Aut(A_5)$ injects into $S_5$; conjugation by all 120
permutations of five letters gives 120 distinct automorphisms, proving
equality. Direct enumeration yields
\[
 |\Epi(F_2,A_5)|=2280,\qquad |\Epi(F_3,A_5)|=200160,
 \qquad h_3(A_5)=1668.
\]
The finite checker is a replay of these small counts, not a proof of
the general group-theoretic statements. The threshold nineteen is
also explicitly present in the primary source \cite{bc}.

\paragraph{Comparison with raw counts.}
If $H$ is $d$-generated then $|\Hom(H,A_5)|\le60^d$. This remains a
valid bound, but does not use independence of simple quotient kernels.
When all nontrivial maps are onto, $|\Hom(H,A_5)|=1+120\kappa_{A_5}(H)$;
the raw count must then reach at least thirty kernels to exceed 3600,
while the joint test needs twenty. Proper-image maps can change this
comparison, so neither test is claimed to dominate the other for every
group. For example, the abstract group $A_5^{20}$ has just 2401 maps
to $A_5$ but is three-generated: all nontrivial maps select one factor
and an automorphism, and $20\le h_3(A_5)$. This is an abstract control,
not an asserted branched-cover group.

\section{A strict improvement on a known knot}

Let $K=P(2,3,3,5)$, defined as the numerator closure of the sum of
rational tangles $1/2+1/3+1/3+1/5$. It is a knot, since exactly one
pretzel parameter is even. Its standard four-tangle diagram has a
four-bridge upper bound; an explicit four-seed Wirtinger propagation
certificate is also retained.

The double cover of a Montesinos link is Seifert fibred. In the
unnormalized conventions of \cite[Section 7.3, p.~2671]{ps}, its group
has the presentation
\[
 H=\langle x_1,\ldots,x_4,h\mid h\text{ central},\
 x_i^{a_i}=h^{-b_i},\ x_1x_2x_3x_4=1\rangle.
\]
An integer twist can be absorbed into the unnormalized pairs; after
killing $h$ it does not change the following quotient:
\[
 \Gamma(a_1,a_2,a_3,a_4)=
 \langle x_1,\ldots,x_4\mid x_i^{a_i}=1,\ x_1x_2x_3x_4=1\rangle.
\]
Every epimorphism $H\to A_5$ kills $h$, since its image is central in
the centerless target. Hence epimorphisms of $H$ and $\Gamma$ to
$A_5$ correspond, with kernels preserved under pullback.

For $(a_1,a_2,a_3,a_4)=(2,3,3,5)$, the exact enumeration of all 1668
generating-triple automorphism classes gives \textbf{23} epimorphism
kernels. Take $x_4=(x_1x_2x_3)^{-1}$ and retain just those triples with
$x_i^{a_i}=1$. The direct knot-diagram computation gives the same
23 kernels, all on one fixed normalized core presentation.

For this pretzel all $b_i=1$. In fact every homomorphism $H\to A_5$
kills $h$: a nonidentity image has prime order $p\in\{2,3,5\}$, and
the relation with exponent $a_i=p$ would require $x^p=h^{-1}$.
No element of $A_5$, whose orders are $1,2,3,5$, can satisfy this.
The full homomorphism count, by image order, is
\[
\begin{array}{c|rrrr|r}
 |\operatorname{image}|&1&3&12&60&\text{total}\\\hline
 \text{number of maps}&1&20&120&2760&2901.
\end{array}
\]
Thus $2760/120=23$ kernels. The raw estimate only forces
$\rk H\ge2$, giving $\mr(K)\ge3$. Corollary~\ref{cor:twenty} instead
gives $\mr(K)\ge4$, matching the certified upper bound. This knot is
already covered by the Montesinos theorem of Boileau--Zieschang
(1985), restated in \cite[Theorem 1]{musick}. It is a positive control
for the stronger certificate, not a new resolution.

The same argument gives a uniform known family: if a four-tangle
Montesinos knot has denominators divisible, after reordering, by
$2,3,3,5$, its cover group surjects onto $\Gamma(2,3,3,5)$.
The 23 distinct kernels pull back to 23 distinct kernels of that one
cover group. Its standard four-bridge upper bound therefore matches
the new lower bound. Numerators and integer twist do not affect this
epimorphism argument.

\section{A finite denominator classification and the candidate panel}

For $A_5$, the order constraints depend only on
$\gcd(a_i,30)$. With one even denominator and three odd denominators
each divisible by three or five, all forty reduced patterns give
the following exact kernel counts. Rows list the three odd gcds;
columns give the even gcd. Each entry concerns its own single group
$\Gamma$; entries are never added together.
\begin{center}
\begin{tabular}{lrrrr}
\toprule
Odd gcds & $2$ & $6$ & $10$ & $30$\\
\midrule
$(3,3,3)$&12&21&40&49\\
$(3,3,5)$&23&50&62&89\\
$(3,3,15)$&35&71&101&137\\
$(3,5,5)$&29&67&70&108\\
$(3,5,15)$&51&115&129&193\\
$(3,15,15)$&85&184&226&325\\
$(5,5,5)$&32&74&82&124\\
$(5,5,15)$&60&138&150&228\\
$(5,15,15)$&109&248&274&413\\
$(15,15,15)$&191&425&491&725\\
\bottomrule
\end{tabular}
\end{center}
Thirty-nine patterns reach twenty. The exception $(2,3,3,3)$ has
twelve, so this $A_5$ test is not universal even among known
four-tangle pretzel controls. If an odd denominator is divisible by
neither three nor five, its cone generator must map to the identity;
the remaining three-cone group is two-generated, giving at most
nineteen kernels. Four-odd-denominator patterns are outside this
finite panel.

\paragraph{All 293 retained uncertain names.}
The second panel was frozen before testing, comprising all 130
alternating and 163 nonalternating thirteen-crossing names left in the
pinned October release \cite{deposit}, with unchanged order and no
target replacement. PD data and reported $[3,4]$ status come from the
pinned KnotInfo 2026.11.2 selection. Every diagram has a checked
four-seed upper certificate. The complete per-knot distribution is
\[
\begin{array}{c|rrrrr}
 \kappa_{A_5}(\pi_1\Sigma_2K)&0&1&2&3&5\\\hline
 \text{number of knots}&151&118&22&1&1.
\end{array}
\]
The maximum five occurs at $13n\_4587$; $13n\_4589$ has three.
None reaches twenty. For the first sixteen previously tested targets,
the counts in retained order are
\[
 1,2,1,2,1,2,1,2,0,1,0,1,0,1,0,1.
\]
Their old raw counts are exactly $1+120\kappa$, consistent with the
earlier independent replay. The twelve first maps on different knots
remain twelve separate witnesses and cannot certify a joint rank gain.

\section{Enumeration completeness, replay and remaining work}

The normalized core presentation fixes one arc generator to the
identity. At each crossing its generators obey
$z=x y^{-1}x$, with $x$ the overarc. Killing the normalized arc
identifies this group with $\pi_1(\Sigma_2K)$ by the classical
Kelly--Wada core-group construction \cite{core}. A four-seed
propagation trace expresses every remaining arc in three generators.
For an epimorphism those three images generate $A_5$. Testing all
1668 generating-triple automorphism classes is therefore exhaustive
for quotient kernels, irrespective of possible proper-image maps.
Full crossing relations are checked after propagation. Each retained
epimorphism carries the identifier of its one normalized source
presentation; automorphic duplicates and cross-source pooling are
explicitly rejected by negative controls.

The producer uses even permutations and earliest-nonidentity-entry
normalization of tuples. The separate verifier constructs
$\operatorname{PSL}_2(\mathbb F_5)$ and its 120 conjugation
automorphisms from $\operatorname{PGL}_2(\mathbb F_5)$, then marks
all tuple orbits exhaustively. Its iterative crossing closure does
not use the producer's fixed propagation trace. It reproduces all
293 counts and all forty orbifold counts. It directly checks every
retained permutation epimorphism and its automorphism class.
The PD parser is shared; this is disclosed rather than described as
wholly independent software.

The known pretzel's 23 kernels and full 2901-map count are replayed
both from its diagram and from the cone presentation, with image
orders checked. Two separate rational-tangle constructors, one using
Spherogram and one paired-endpoint graphs, agree by an exact
crossing-preserving planar-map bijection. A connected sum of three
trefoils supplies another known control: its cover group is
$C_3*C_3*C_3$ and has 49 kernels; two trefoils give only one. The
free-group controls recover nineteen kernels for $F_2$, 1668 for
$F_3$, and an explicit list of twenty distinct maps on one $F_3$.

\clearpage
The receipt \path{joint_a5_verification.json} seals the finite inputs
and reports the replay. These checks do not replace the uniform
proofs, which are unformalized and externally unreviewed. The
candidate panel supplies no new sharp case. A further campaign needs
a structural reason to obtain more independent kernels for one
group, another finite target with an appropriate capacity threshold,
or a different cover-rank certificate. Repeating the completed
$A_5$ panel or pooling witnesses across knots cannot improve it.

\begin{thebibliography}{9}
\bibitem{bc} Burness--Crestani, On the generating graph of direct powers
of a simple group, J. Algebraic Combinatorics 38 (2013), 329--350.
\url{https://doi.org/10.1007/s10801-012-0405-x}.
Author manuscript, p.~2:
\url{https://seis.bristol.ac.uk/~tb13602/docs/BC_1.pdf}.
Their reference [18] is Hall, The Eulerian functions of a group,
Quart. J. Math. Oxford 7 (1936), 134--151; the original Hall paper
has not been independently read here.
\bibitem{ps} Poudel--Saveliev, Link homology and equivariant gauge
theory, Algebraic \& Geometric Topology 17 (2017), 2635--2685,
Section 7.3. \url{https://doi.org/10.2140/agt.2017.17.2635}.
\bibitem{musick} Musick, Minimal bridge projections for 11-crossing
prime knots (2012), Theorem 1, restating Boileau--Zieschang (1985).
\url{https://arxiv.org/abs/1208.4233v3}.
\bibitem{core} Silver--Traldi--Williams, Core groups (2024),
introduction and Section 7. \url{https://arxiv.org/abs/2311.02048v2}.
\bibitem{deposit} Bhattacharjee, unrefereed release 1.3.0 (2026),
pinned uncertainty list. \url{https://doi.org/10.5281/zenodo.23561157}.
\end{thebibliography}
\end{document}
